\documentclass[11pt,a4paper]{article}
\usepackage[margin=2.5cm]{geometry}
\usepackage{booktabs}
\usepackage{amsmath,amssymb}
\usepackage{multirow}
\usepackage{graphicx}
\usepackage{xcolor}
\usepackage{hyperref}
\usepackage{natbib}
\usepackage{subcaption}
\usepackage{array}

\newcommand{\sica}{\textsc{Sica}}
\newcommand{\effk}{\text{Eff-K}}
\newcommand{\pp}{\text{pp}}

\title{When Do Self-Extracted Symbolic Constraints Help LLM Reasoning?\\A Diagnostic Study}

\author{Anonymous}
\date{}

\begin{document}
\maketitle

\begin{abstract}
Self-consistency (SC) improves large language model (LLM) reasoning by sampling multiple chains of thought and taking a majority vote.
A natural extension is to extract formal constraints from reasoning traces and aggregate them with a MAX-SAT solver, selecting answers that satisfy the most constraints.
We conduct a systematic diagnostic study of this approach, which we call Symbolic Invariant Consensus Aggregation (\sica{}), across 7 models and 6 reasoning domains.
Our oracle control establishes that MAX-SAT aggregation over gold first-order logic (FOL) annotations yields +14.2 percentage points over SC ($p < 10^{-6}$), confirming that the aggregation mechanism itself is sound.
However, when models extract their own constraints, \sica{} produces zero or negative gains in 9 of 12 tested conditions.
The single reliably positive condition (Mistral-7B on FOLIO, $\Delta$=+4.9\pp{}, $p$=0.013) does not generalize across seeds (1 of 3 seeds yields a null result).
We characterize three failure regimes, quantify trace non-independence via Fleiss' $\kappa$ (effective sample sizes of 1.2--3.6 from nominal $K$=12), and show that 9 targeted remediation strategies all fail to close the oracle gap.
We propose a zero-cost diagnostic based on the multi-candidate rate (C0) and SC accuracy that identifies conditions where \sica{} is guaranteed to be ineffective.
\end{abstract}

\section{Introduction}

Self-consistency \citep{wang2023selfconsistency} is one of the most widely adopted strategies for improving LLM reasoning at inference time.
By sampling $K$ reasoning traces and selecting the most frequent answer via majority vote, SC exploits the observation that correct reasoning paths tend to converge.
This approach has been extended to more sophisticated aggregation: extracting formal constraints from traces and using satisfiability solvers to select answers that are maximally consistent with the extracted knowledge \citep{pan2023logiclm,olausson2023linc}.

The theoretical foundation for such voting-based aggregation is the Condorcet Jury Theorem \citep{condorcet1785}, which guarantees that majority voting converges to the correct answer as the number of independent voters grows, provided each voter is correct with probability $p > 0.5$.
However, this guarantee depends critically on voter independence.
When LLM reasoning traces are sampled from the same model, the independence assumption is severely violated: traces share the same parametric knowledge, the same systematic biases, and often converge to the same errors.

In this work, we ask: \emph{under what conditions can self-extracted symbolic constraints improve over simple majority voting?}
We conduct a comprehensive diagnostic study of Symbolic Invariant Consensus Aggregation (\sica{}), a pipeline that extracts first-order logic constraints from sampled traces and aggregates them via weighted MAX-SAT solving.
Our study spans 7 model families (Qwen-7B/14B/72B, Mistral-7B, LLaMA-3.1-8B, DeepSeek-R1-7B, QwQ-32B) across 6 reasoning domains (FOLIO, ProofWriter, LogiQA, StrategyQA, GSM8K, ARC-Challenge).

Our contributions are:
\begin{enumerate}
\item \textbf{Oracle gap localization.} Gold FOL constraints with MAX-SAT yield +14.2\pp{} over SC ($p < 10^{-6}$), proving the aggregation mechanism is sound. Self-extracted constraints capture only 35\% of this theoretical potential, localizing the bottleneck to extraction quality.
\item \textbf{Systematic failure characterization.} We identify three failure regimes (high-accuracy confirmation bias, low-accuracy noise, domain non-formalizability) and document that \sica{} yields zero or negative gains in 9 of 12 conditions. The single positive condition is fragile across seeds and hyperparameters.
\item \textbf{Remediation breadth.} Nine targeted interventions spanning prompting, architecture, capability, verification, diversity, and filtering all fail to close the oracle gap, suggesting that the self-extraction bottleneck is structural.
\item \textbf{Deployment diagnostic.} The multi-candidate rate (C0) combined with SC accuracy provides a zero-cost screening tool: C0 $< 5\%$ guarantees \sica{} $\equiv$ SC, eliminating unnecessary computation in the majority of deployment scenarios.
\end{enumerate}

\section{Related Work}

\paragraph{Self-Consistency and Extensions.}
Wang et al.~\citeyearpar{wang2023selfconsistency} introduced SC for chain-of-thought reasoning.
Subsequent work has explored universal SC \citep{chen2023usc}, diversity-aware sampling \citep{li2023diverse}, weighted voting \citep{taubenfeld2025weighted}, and scaling analysis of SC \citep{loo2025sc}.
These approaches share the fundamental assumption that sampling diverse traces provides complementary information, an assumption we directly test.

\paragraph{Neuro-Symbolic Reasoning.}
Several systems combine LLM reasoning with formal logic.
Logic-LM \citep{pan2023logiclm} translates problems to symbolic form and uses solvers for inference.
LINC \citep{olausson2023linc} performs neurosymbolic reasoning via in-context learning.
More recent work includes DyLoS \citep{dylos2026}, which trains dedicated extractors, Aristotle \citep{aristotle2025}, which uses structured reasoning with external verifiers, and FoVer \citep{fover2025}, which performs formal verification of LLM outputs.
Our work differs from these in studying \emph{self-extraction} without trained extractors or external verifiers.

\paragraph{Condorcet Jury Theorem and LLM Ensembles.}
The connection between CJT and LLM voting has been explored by \citet{li2024independence}, who study independence violations in cross-model sentiment analysis ensembles.
Our work applies CJT analysis to \emph{within-model} trace sampling, quantifying independence via Fleiss' $\kappa$ \citep{fleiss1971} and introducing the effective sample size metric $\effk{}$.

\paragraph{Test-Time Compute and Verification.}
Scaling inference-time computation through repeated sampling \citep{snell2025scaling}, process reward models \citep{lightman2023prm}, outcome reward models \citep{cobbe2021orm}, and extended reasoning \citep{deepseekr1} represent alternatives to symbolic aggregation.
Our oracle gap analysis clarifies when symbolic approaches have theoretical headroom over these alternatives.

\section{Method: \sica{} Pipeline}

\subsection{Overview}

Given a reasoning problem $q$ and an LLM $\mathcal{M}$, \sica{} proceeds in three stages:

\paragraph{Stage 1: Trace Sampling.}
Sample $K$ reasoning traces $\{t_1, \ldots, t_K\}$ from $\mathcal{M}$ at temperature $T$, each producing a candidate answer $a_i$.

\paragraph{Stage 2: Constraint Extraction.}
For each trace $t_i$, prompt $\mathcal{M}$ to extract first-order logic (FOL) constraints $\{c_{i,1}, \ldots, c_{i,m_i}\}$ that formalize the reasoning steps. Each constraint is a FOL formula over the candidate answer set.

\paragraph{Stage 3: MAX-SAT Aggregation.}
Encode each constraint as a weighted clause in a MAX-SAT instance. For each candidate answer $a$, compute the total weight of satisfied constraints. Select the answer maximizing total satisfied weight:
\begin{equation}
a^* = \arg\max_{a \in \mathcal{A}} \sum_{i=1}^{K} \sum_{j=1}^{m_i} w_{i,j} \cdot \mathbf{1}[c_{i,j} \text{ satisfied by } a]
\end{equation}

\subsection{Independence Quantification}

We measure trace agreement using Fleiss' $\kappa$ \citep{fleiss1971} computed over the $K$ traces' answer distributions per problem. The effective number of independent voters is:
\begin{equation}
\effk{} = \frac{K}{1 + (K-1)\kappa}
\end{equation}
where $\kappa$ is the average pairwise agreement. When $\kappa = 0$ (full independence), $\effk{} = K$. When $\kappa = 1$ (perfect agreement), $\effk{} = 1$.

\subsection{Diagnostic Metrics}

\paragraph{Multi-Candidate Rate (C0).}
The fraction of problems where at least two distinct candidate answers appear across $K$ traces. Problems with C0 = 0 (unanimous agreement) are trivially identical under SC and \sica{}.

\paragraph{Oracle Control.}
To separate extraction quality from aggregation quality, we replace self-extracted constraints with gold-standard FOL annotations and measure the resulting MAX-SAT accuracy.

\section{Experimental Setup}

\paragraph{Models.}
We evaluate 7 models spanning 5 families: Qwen2.5-7B, Qwen2.5-14B, Qwen2.5-72B-AWQ (4-bit), LLaMA-3.1-8B, Mistral-7B-Instruct-v0.3, DeepSeek-R1-Distill-Qwen-7B, and QwQ-32B.

\paragraph{Datasets.}
FOLIO \citep{han2022folio} (natural language FOL inference, 204 problems), ProofWriter \citep{tafjord2021proofwriter} (synthetic deductive reasoning), LogiQA \citep{liu2023logiqa} (logical comprehension), StrategyQA \citep{geva2021strategyqa} (multi-hop commonsense, 60 problems), ARC-Challenge \citep{clark2018arc} (science QA, 60 problems), and GSM8K \citep{cobbe2021gsm8k} (math word problems, 30 problems). D012 denotes the combined FOLIO+LogiQA+ProofWriter subset (60 problems).

\paragraph{Configuration.}
Default: $K$=12, $T$=0.7, seed 42. Sensitivity analyses vary $K \in \{4, 6, 8, 11, 12, 16, 20\}$ and $T \in \{0.3, 0.5, 0.7, 1.0\}$. Statistical significance is assessed via McNemar's test.

\section{Results}

\subsection{Main Results}

Table~\ref{tab:main} presents \sica{} versus SC across all model-domain conditions.
Of 12 primary conditions, only one (Mistral-7B on FOLIO-204) yields a statistically significant positive result ($\Delta$=+4.9\pp{}, $p$=0.013).
Three conditions show positive but non-significant trends (LLaMA-8B on D012, +6.7\pp{}, $p$=0.219; Qwen-14B on StrategyQA, +3.3\pp{}, $p > 0.1$; QwQ-32B on AIME, $\Delta$=0.0).
The remaining conditions show zero or negative deltas.

\begin{table}[t]
\centering
\caption{Main results: SC vs.\ \sica{} across models and domains. $n$: number of problems. C0: multi-candidate rate. $\Delta$: SICA $-$ SC in percentage points. Bold: statistically significant ($p < 0.05$).}
\label{tab:main}
\small
\begin{tabular}{llrrrrrr}
\toprule
Model & Domain & $n$ & SC\% & \sica{}\% & $\Delta$\pp{} & $p$ & C0 \\
\midrule
Mistral-7B & FOLIO-204 & 204 & 54.4 & 59.3 & \textbf{+4.90} & \textbf{0.013} & 0.69 \\
Qwen2.5-14B & FOLIO & 204 & 75.5 & 75.7 & +0.25 & n.s. & 0.38 \\
Qwen2.5-7B & D012 & 60 & 68.3 & 68.3 & 0.00 & 1.000 & 0.43 \\
Qwen2.5-72B & D012 & 60 & 83.3 & 81.7 & $-$1.67 & 0.625 & 0.20 \\
LLaMA-3.1-8B & D012 & 60 & 60.0 & 66.7 & +6.67 & 0.219 & 0.62 \\
DeepSeek-R1-7B & D012 & 60 & 71.7 & 68.3 & $-$3.33 & 0.480 & 0.58 \\
QwQ-32B & AIME & 29 & --- & --- & 0.00 & n.s. & 0.03 \\
QwQ-32B & PW & 17 & --- & --- & $-$5.90 & n.s. & 0.03 \\
\midrule
Qwen2.5-14B & StrategyQA & 60 & 53.3 & 56.7 & +3.33 & n.s. & 0.10 \\
Qwen2.5-14B & ARC-Chall. & 60 & 96.7 & 95.0 & $-$1.67 & --- & 0.05 \\
Qwen2.5-14B & GSM8K & 30 & 96.7 & 96.7 & 0.00 & --- & 0.13 \\
\bottomrule
\end{tabular}
\end{table}

Several patterns emerge.
The Qwen model family shows zero or negative gains at all scales (7B, 14B, 72B), consistent with high baseline accuracy and low C0 values.
LLaMA-8B shows a positive trend (+6.7\pp{}) on D012 ($n$=60), but this result has $p$=0.219 and is accompanied by 44 truncated traces (of 60 problems) and an answer extraction rate of only 89.6\%, raising concerns about pipeline artifacts.
QwQ-32B, a reasoning-specialized model, has C0 $< 5\%$ across both domains (aggregate C0 effective = 2.6\%), confirming that near-deterministic trace generation renders \sica{} equivalent to SC.
GSM8K exhibits a domain mismatch: FOL extraction succeeds on only 14.7\% of traces, as mathematical reasoning requires numerical constraints rather than propositional logic.

\subsection{Oracle Gap Analysis}

Table~\ref{tab:oracle} presents the oracle control experiment on Qwen2.5-14B FOLIO.
Gold FOL annotations with MAX-SAT aggregation yield 89.7\% accuracy, a +14.2\pp{} improvement over SC ($p = 2.0 \times 10^{-7}$).
Self-extracted \sica{} achieves only +0.25\pp{}, capturing approximately 2\% of the oracle potential.
Even in the best self-extraction condition (Mistral-7B FOLIO, +4.9\pp{}), the utilization rate is only 35\%.

\begin{table}[t]
\centering
\caption{Oracle control on Qwen2.5-14B FOLIO ($n$=204). Gold FOL: human-annotated first-order logic constraints. Self: model-extracted constraints.}
\label{tab:oracle}
\small
\begin{tabular}{lrrr}
\toprule
Method & Accuracy\% & $\Delta$ vs.\ SC & $p$-value \\
\midrule
SC (majority vote) & 75.5 & --- & --- \\
\sica{} (self-extracted) & 75.7 & +0.25\pp{} & n.s. \\
\sica{} (gold FOL) & 89.7 & +14.22\pp{} & $2.0 \times 10^{-7}$ \\
\bottomrule
\end{tabular}
\end{table}

This result has two implications.
First, the MAX-SAT aggregation mechanism is not the bottleneck; when given high-quality constraints, it substantially outperforms majority voting.
Second, the gap between oracle and self-extraction performance localizes the failure to constraint quality.
A two-layer analysis reveals the mechanism: 43\% of self-extracted constraints are trivially satisfied by all candidate answers (providing no discriminative signal), and 78\% of constraints mirror the model's own answer (confirmation bias).

\subsection{Independence Quantification}

Table~\ref{tab:independence} reports Fleiss' $\kappa$ and $\effk{}$ across model-domain conditions.
For $K$=12 nominal traces, the effective number of independent voters ranges from 1.20 (Qwen-14B LogiQA) to 3.63 (Mistral-7B ProofWriter-600).

\begin{table}[t]
\centering
\caption{Trace independence metrics. $\kappa$: Fleiss' kappa. $\effk{}$: effective independent voters for $K$=12. SC: self-consistency accuracy. $\Delta$: \sica{}$-$SC.}
\label{tab:independence}
\small
\begin{tabular}{llrrrr}
\toprule
Model & Domain (seed) & $\kappa$ & $\effk{}$ & SC\% & $\Delta$\pp{} \\
\midrule
Qwen-14B & FOLIO & 0.744 & 1.31 & 75.5 & +0.25 \\
Qwen-14B & LogiQA & 0.821 & 1.20 & 78.5 & +0.50 \\
Mistral-7B & FOLIO (s42) & 0.500 & 1.85 & 54.4 & +4.90 \\
Mistral-7B & FOLIO (s123) & 0.412 & 2.17 & 54.4 & +3.43 \\
Mistral-7B & FOLIO (s456) & 0.401 & 2.22 & 54.4 & $-$0.49 \\
Mistral-7B & PW-600 & 0.210 & 3.63 & 39.3 & $-$2.83 \\
\bottomrule
\end{tabular}
\end{table}

A cross-model pattern is clear: Qwen-14B ($\effk{} \approx 1.2$--$1.3$) consistently shows zero gain, while Mistral-7B ($\effk{} \approx 1.9$--$2.2$ on FOLIO) shows variable gains.
However, $\effk{}$ has no within-model predictive power.
Among the three Mistral-7B FOLIO seeds, seed 42 has the \emph{lowest} $\effk{}$ (1.85) yet the highest $\Delta$ (+4.9\pp{}), while seed 456 has higher $\effk{}$ (2.22) but negative $\Delta$ ($-$0.49\pp{}).
Mistral-7B on ProofWriter-600 has the highest $\effk{}$ of all conditions (3.63) yet produces a significant \emph{negative} $\Delta$ ($-$2.83\pp{}, $p$=0.010), because SC accuracy (39.3\%) falls below the quality threshold.

To quantify this predictive failure, we evaluated joint $\effk{}$+SC prediction across all 12 conditions (Table~\ref{tab:prediction}).
The framework correctly predicts the sign of $\Delta$ in only 6 of 12 conditions (50\%), equivalent to chance.
Key failures include Mistral-7B PW-600 (highest $\effk{}$, predicted positive, actual $-$2.83\pp{}) and within-seed reversals where similar $\effk{}$+SC values produce opposite-sign deltas.

\begin{table}[t]
\centering
\caption{Joint $\effk{}$+SC prediction accuracy across 12 conditions. The framework predicts the sign of $\Delta$ at chance level (50\%).}
\label{tab:prediction}
\small
\begin{tabular}{lrr}
\toprule
Prediction outcome & Count & Fraction \\
\midrule
Correctly predicted positive & 2 & 16.7\% \\
Correctly predicted negative/zero & 4 & 33.3\% \\
Incorrectly predicted & 6 & 50.0\% \\
\midrule
Total & 12 & 100\% \\
\bottomrule
\end{tabular}
\end{table}

\subsection{K Sensitivity}

We vary $K$ for both Mistral-7B (FOLIO-204) and Qwen2.5-14B (D012, 60 problems). Results appear in Table~\ref{tab:ksens}.

\begin{table}[t]
\centering
\caption{$K$ sensitivity. Left: Mistral-7B on FOLIO-204 (seed 42). Right: Qwen2.5-14B on D012 (60 problems). None of the pairwise differences are statistically significant.}
\label{tab:ksens}
\small
\begin{tabular}{crrrcrr}
\toprule
& \multicolumn{3}{c}{Mistral-7B / FOLIO-204} & & \multicolumn{2}{c}{Qwen-14B / D012} \\
\cmidrule{2-4} \cmidrule{6-7}
$K$ & SC\% & \sica{}\% & $\Delta$\pp{} & & $K$ & $\Delta$\pp{} \\
\midrule
4 & 54.4 & 53.4 & $-$0.98 & & 6 & 0.0 \\
8 & 53.9 & 55.9 & +1.96 & & 8 & $-$1.7 \\
12 & 54.4 & 59.3 & +4.90 & & 11 & 0.0 \\
16 & 53.9 & 55.4 & +1.47 & & 12 & $-$5.0 \\
20 & 54.9 & 52.4 & $-$2.45 & & 16 & $-$3.3 \\
 & & & & & 20 & $-$3.3 \\
\bottomrule
\end{tabular}
\end{table}

For Mistral-7B, $K$=12 produces the largest $\Delta$ (+4.9\pp{}), but the pattern across $K$ values does not follow a smooth trend.
The sequence ($-$1.0, +2.0, +4.9, +1.5, $-$2.5) is consistent with random fluctuation at $n$=204; confidence intervals overlap substantially, and we cannot reject the null hypothesis that all $K$ values produce the same $\Delta$.
For Qwen-14B, all $K$ values yield zero or negative $\Delta$, with monotonic degradation at $K \geq 12$, consistent with redundant correlated constraints accumulating noise.

\subsection{Temperature Sensitivity}

Table~\ref{tab:tsens} reports \sica{} performance at four temperature values and a multi-temperature mixing condition for Mistral-7B on FOLIO-204.

\begin{table}[t]
\centering
\caption{Temperature sensitivity for Mistral-7B on FOLIO-204 ($K$=12). Multi-temp mixes equal numbers of traces at $T$=0.3, 0.7, and 1.0. $\effk{}$ computed from Fleiss' $\kappa$ at each temperature.}
\label{tab:tsens}
\small
\begin{tabular}{lrrrrl}
\toprule
$T$ & SC\% & \sica{}\% & $\Delta$\pp{} & $\effk{}$ & $p$-value \\
\midrule
0.3 & 53.4 & 57.4 & +3.92 & 1.45 & 0.065 \\
0.5 & 53.9 & 56.4 & +2.45 & 1.65 & 0.230 \\
0.7 & 54.4 & 59.3 & +4.90 & 1.85 & 0.013 \\
1.0 & --- & --- & +0.49 & --- & 0.739 \\
\midrule
Multi & 54.4 & 56.4 & +2.45 & 2.80 & 0.180 \\
\bottomrule
\end{tabular}
\end{table}

\sica{} shows positive $\Delta$ at $T$=0.3 through $T$=0.7, with only $T$=1.0 collapsing to near zero.
This pattern complicates a simple fragility narrative: the effect is not confined to a narrow temperature window but persists across low-to-medium temperatures, although none of the individual conditions except $T$=0.7 reach significance.
The non-monotonic pattern across temperatures ($T$=0.3: +3.9, $T$=0.5: +2.5, $T$=0.7: +4.9, $T$=1.0: +0.5) lacks a clean functional form.

The multi-temperature condition is informative.
Mixing traces across temperatures increases $\effk{}$ from 1.85 (single $T$=0.7) to 2.80, a 51\% improvement in effective independence.
The resulting $\Delta$ of +2.45\pp{} ($p$=0.18), while non-significant, is positive, and this is the only remediation strategy that increases both diversity \emph{and} shows a positive trend.
This result weakens a blanket claim that diversity interventions cannot help, though the evidence remains underpowered.

\subsection{Constraint Quality Analysis}

\paragraph{Qualitative analysis.}
We analyze constraint quality on 8 representative problems.
Of these, 45\% contain only trivially satisfied constraints (satisfied by all candidate answers), 6 of 8 exhibit confirmation bias (constraints mirror the model's own answer), and only 1 of 8 contains a genuinely informative constraint.
Overall, 78\% of extracted constraints agree with the model's majority-vote answer, indicating that the extraction process largely recapitulates the model's existing beliefs rather than introducing independent evidence.

\paragraph{Quantitative correlation.}
Across 10 model-domain conditions, the correlation between constraint precision and $\Delta$ is $r$=0.28 ($p$=0.43, $n$=10), far from significant.
Similarly, C0 alone correlates with $\Delta$ at $r$=0.31 ($p$=0.33, $n$=12), and correctly predicts the direction of $\Delta$ in only 5 of 12 conditions (41.7\%).
High C0 is associated with both positive deltas (LLaMA-8B D012: C0=0.62, $\Delta$=+6.7\pp{}) and negative deltas (DeepSeek-R1 D012: C0=0.58, $\Delta$=$-$3.3\pp{}), indicating that answer diversity is necessary but not sufficient for \sica{} effectiveness.

\paragraph{Aggregation mechanism ablation.}
To test whether MAX-SAT solving contributes beyond simple constraint counting, we compare three aggregation strategies on 804 problems from the primary dataset (Table~\ref{tab:ablation}).
MAX-SAT and count-only agree on 803 of 804 problems.
No method significantly outperforms the SC baseline.
This result indicates that the symbolic optimization machinery (weighted MAX-SAT) provides negligible value over simple constraint counting, because constraint satisfaction scores are nearly uniform across candidates (mean satisfaction ratio $\approx$ 0.96).

\begin{table}[t]
\centering
\caption{Aggregation mechanism ablation ($n$=804). Agreement: fraction of problems where method and count-only produce the same answer.}
\label{tab:ablation}
\small
\begin{tabular}{lrr}
\toprule
Method & Accuracy\% & Agreement w/ count-only \\
\midrule
SC (majority vote) & 76.2 & --- \\
MAX-SAT (weighted) & 76.5 & --- \\
Count-only & 76.6 & 803/804 \\
Random-weight & 75.9 & --- \\
\bottomrule
\end{tabular}
\end{table}

\subsection{Seed Replication}

The primary positive result (Mistral-7B FOLIO-204, $K$=12, $T$=0.7) was replicated across three random seeds.
Two of three seeds yield significant positive results: seed 42 ($\Delta$=+4.90\pp{}, $p$=0.013) and seed 123 ($\Delta$=+3.43\pp{}, $p$=0.039).
Seed 456 produces a null result ($\Delta$=$-$0.49\pp{}, n.s.).
The effect is real but fragile: even under identical hyperparameters and the same model-domain combination, random seed variation can eliminate or reverse the gain.

\subsection{Remediation Attempts}

Table~\ref{tab:remediation} summarizes 9 remediation strategies targeting different hypothesized bottlenecks.

\begin{table}[t]
\centering
\caption{Nine remediation strategies and their outcomes. All were applied to Mistral-7B FOLIO. ``Outcome'' indicates the effect on $\Delta$ relative to the base condition.}
\label{tab:remediation}
\small
\begin{tabular}{p{3.5cm}p{3cm}p{4.5cm}}
\toprule
Strategy & Target & Outcome \\
\midrule
Contrastive prompting & Extraction quality & No improvement \\
Cross-model extraction & Independence & No improvement \\
Generational scaling & Model capability & No improvement \\
Solver-guided refinement & Constraint consistency & No improvement \\
Consistency filtering & Noise reduction & No improvement \\
Cross-trace UNSAT & Contradiction detection & No improvement \\
Discriminative filtering & Non-informative removal & 0 gains (92\% already discriminative) \\
$T$=1.0 diversity & Trace diversity & $\Delta$ collapses: +4.9 $\to$ +0.5\pp{} \\
Multi-temperature & Structured diversity & +2.45\pp{} (n.s., $p$=0.18) \\
\bottomrule
\end{tabular}
\end{table}

Eight of nine strategies produce no improvement or degradation.
The single partial exception is multi-temperature mixing, which increases $\effk{}$ from 1.85 to 2.80 and yields $\Delta$=+2.45\pp{}, positive but non-significant.
This result suggests that structured diversity interventions may have marginal value, but the effect size is smaller than the single-temperature optimum (+4.9\pp{} at $T$=0.7) and not statistically reliable.
Notably, simply increasing temperature to $T$=1.0 collapses $\Delta$ from +4.9 to +0.5\pp{}, demonstrating that unstructured diversity is harmful.

The discriminative filtering result is particularly informative: 92\% of constraints are already discriminative (they assign different scores to different candidates), yet removing the 8\% non-discriminative constraints produces zero improvement.
The problem is not that constraints lack discriminative power, but that their discriminative direction aligns with the model's existing biases.

\subsection{Additional Baselines}

\paragraph{Weighted SC.}
We test three weighting schemes for SC: inverse-length, length-weighted, and constraint-count weighted (Table~\ref{tab:weightedsc}).
No scheme significantly outperforms uniform SC, consistent with the thesis that re-weighting correlated votes does not address the underlying independence problem.

\begin{table}[t]
\centering
\caption{Weighted SC variants on the primary dataset ($n$=804). All differences are non-significant.}
\label{tab:weightedsc}
\small
\begin{tabular}{lr}
\toprule
Weighting & Accuracy\% \\
\midrule
Uniform (SC) & 76.2 \\
Inverse-length & 76.7 \\
Length-weighted & 75.3 \\
Constraint-count & 75.3 \\
\bottomrule
\end{tabular}
\end{table}

\paragraph{Logic-LM baseline.}
We compare against a Logic-LM-style approach that performs single-trace FOL verification across $K$ traces and takes a majority vote.
On FOLIO, this achieves 75.5\% (vs.\ SC 75.0\%, \sica{} 74.5\%).
All three methods are statistically indistinguishable, suggesting that multi-trace constraint aggregation provides no advantage over per-trace symbolic verification when traces are highly correlated.

\subsection{Deployment Diagnostic: C0 + SC}

We propose a two-dimensional diagnostic for screening \sica{} applicability before deployment.

\paragraph{C0 screening.}
When C0 $< 5\%$, \sica{} is mathematically equivalent to SC (no multi-candidate problems exist for MAX-SAT to resolve).
This condition holds for QwQ-32B (C0 = 2.6\%), ARC-Challenge (C0 = 5\%), and near-threshold for GSM8K (C0 = 13.3\%, but with 14.7\% extraction success).
C0 screening eliminates approximately half of model-domain conditions from consideration at zero computational cost.

\paragraph{SC accuracy regime.}
Among conditions where C0 is sufficient, SC accuracy further stratifies outcomes:
\begin{itemize}
\item SC $> 75\%$ (Type A): Confirmation bias dominates. Constraints mirror correct answers, providing no additional signal. Examples: Qwen-14B FOLIO (75.5\%), Qwen-72B D012 (83.3\%).
\item SC $< 40\%$ (Type C): Trace quality is insufficient. Constraints encode mostly incorrect reasoning, degrading performance. Example: Mistral-7B PW-600 (39.3\%, $\Delta$=$-$2.83\pp{}).
\item SC $\approx$ 50--55\% with high C0 (candidate zone): The only observed positive results fall here. Example: Mistral-7B FOLIO (54.4\%, C0=0.69, $\Delta$=+4.9\pp{}).
\end{itemize}

However, we emphasize that this diagnostic has limited positive predictive power.
While C0 $< 5\%$ reliably predicts $\Delta \approx 0$ (necessary condition), high C0 combined with moderate SC does not guarantee positive $\Delta$ (the prediction accuracy is 41.7\% across all conditions).
The diagnostic is useful for ruling out \sica{}, not for recommending it.

\subsection{Overhead Analysis}

\sica{} approximately doubles wall-clock time per problem: mean total time is 61.6 seconds, with trace generation consuming 51.3\% and constraint extraction consuming 47.0\%.
MAX-SAT solving accounts for less than 0.01\% of total time.
Given that $\Delta \leq 0$ in the majority of conditions, the cost-benefit ratio is unfavorable for general deployment.
The overhead analysis does suggest that future work should focus on extraction quality rather than solver efficiency, as the computational bottleneck and the quality bottleneck coincide in the extraction stage.

\section{Discussion}

\paragraph{The self-extraction trap.}
Our results paint a consistent picture: the MAX-SAT aggregation mechanism is effective when given high-quality constraints (oracle gap: +14.2\pp{}), but self-extraction from the same model that generated the reasoning traces produces constraints that largely recapitulate the model's existing beliefs.
This is not a simple quality problem (92\% of constraints are discriminative) but a directionality problem: constraints discriminate in the same direction as the model's majority vote, providing no orthogonal signal.

\paragraph{Limitations of the independence framework.}
While Fleiss' $\kappa$ and $\effk{}$ provide a useful descriptive vocabulary (explaining why Qwen-family models uniformly show zero gain), the framework has no predictive power within a model family.
The 50\% joint prediction accuracy across 12 conditions is equivalent to random guessing.
We therefore present $\effk{}$ as a diagnostic tool for cross-model comparison, not as a predictive framework.

\paragraph{Multi-temperature as a promising direction.}
The multi-temperature result ($\Delta$=+2.45\pp{}, $\effk{}$ increase from 1.85 to 2.80) is the only remediation showing a positive trend.
While underpowered in our study ($p$=0.18), this suggests that structured diversity interventions deserve further investigation, particularly at larger sample sizes.
This partially qualifies our ``9/9 failed'' framing: 8 strategies clearly failed, and 1 showed marginal promise.

\paragraph{Limitations.}
Our primary positive result comes from a single model-domain-temperature combination (Mistral-7B, FOLIO, $T$=0.7).
Sample sizes for pilot domains (StrategyQA, ARC, GSM8K) are small ($n$=30--60), limiting statistical power.
The LLaMA-8B result (+6.7\pp{}) may be confounded by high trace truncation rates (73\%).
We do not evaluate trained extractors or external verifiers, which represent a different paradigm from the self-extraction setting we study.

\section{Conclusion}

We present a diagnostic study of symbolic consensus aggregation across 7 models and 6 domains.
Three findings stand out.
First, the oracle gap (+14.2\pp{}) proves that constraint-based aggregation is a viable mechanism when constraints are reliable, localizing the current failure to extraction quality.
Second, nine remediation strategies fail to close this gap, with only multi-temperature mixing showing marginal promise, suggesting that the self-extraction bottleneck is largely structural.
Third, the C0+SC diagnostic provides a practical screening tool that can eliminate unnecessary \sica{} computation in the majority of deployment scenarios, despite its limited positive predictive power.
The central lesson is that same-model self-extraction produces constraints that mirror rather than complement the model's reasoning, a fundamental limitation that simple diversity or filtering interventions cannot overcome.

\bibliography{references}
\bibliographystyle{plainnat}

\appendix

\section{Tiebreak Analysis}
\label{app:tiebreak}

When SC produces a tie (equal votes for two answers), \sica{} can break the tie using constraint satisfaction scores. On $n$=5 tiebreak problems, \sica{} achieves 80\% accuracy compared to 50\% for random tiebreak, 60\% for alphabetical, and 40\% for first-seen ordering. The sample size is too small for statistical inference.

\section{QwQ-32B Negative Control}
\label{app:qwq}

QwQ-32B, a reasoning-specialized model, produces near-deterministic traces with aggregate C0 effective of 2.6\%. On AIME ($n$=29), $\Delta$=0.0\pp{}. On ProofWriter ($n$=17), $\Delta$=$-$5.9\pp{} (non-significant given the small sample). This confirms that reasoning models with low trace diversity gain nothing from symbolic aggregation.

\section{Full Overhead Breakdown}
\label{app:overhead}

\begin{table}[h]
\centering
\caption{Wall-time breakdown per problem (mean over all conditions).}
\small
\begin{tabular}{lr}
\toprule
Stage & \% of total time \\
\midrule
Trace generation & 51.3\% \\
Constraint extraction & 47.0\% \\
MAX-SAT solving & $<$0.01\% \\
Other & 1.7\% \\
\midrule
Total (mean) & 61.6 seconds \\
\sica{} overhead vs.\ SC-only & +48.7\% \\
\bottomrule
\end{tabular}
\end{table}

\begin{thebibliography}{40}

\bibitem[Wang et al.(2023)]{wang2023selfconsistency}
Wang, X., Wei, J., Schuurmans, D., Le, Q., Chi, E., Narang, S., Chowdhery, A., and Zhou, D.
\newblock Self-consistency improves chain of thought reasoning in language models.
\newblock In \emph{ICLR}, 2023.

\bibitem[Pan et al.(2023)]{pan2023logiclm}
Pan, L., Alber, M., Chen, W., Gashteovski, K., Xie, Y., Feng, Z., and Lawrence, C.
\newblock Logic-LM: Empowering large language models with symbolic solvers for faithful logical reasoning.
\newblock In \emph{Findings of EMNLP}, 2023.

\bibitem[Olausson et al.(2023)]{olausson2023linc}
Olausson, T.~X., Gu, A., Lipkin, B., Zhang, C.~E., Solar-Lezama, A., Tenenbaum, J.~B., and Levy, R.
\newblock LINC: A neurosymbolic approach for logical reasoning by combining language models with first-order logic provers.
\newblock In \emph{EMNLP}, 2023.

\bibitem[Condorcet(1785)]{condorcet1785}
de Condorcet, M.~J.~A.~N.
\newblock \emph{Essai sur l'application de l'analyse \`{a} la probabilit\'{e} des d\'{e}cisions rendues \`{a} la pluralit\'{e} des voix}.
\newblock 1785.

\bibitem[Austen-Smith and Banks(1996)]{austensmith1996}
Austen-Smith, D. and Banks, J.~S.
\newblock Information aggregation, rationality, and the {C}ondorcet jury theorem.
\newblock \emph{American Political Science Review}, 90(1):34--45, 1996.

\bibitem[Li et al.(2024)]{li2024independence}
Li, Z., et al.
\newblock Examining independence in ensemble sentiment classification: A {C}ondorcet jury theorem perspective.
\newblock \emph{arXiv:2409.00094}, 2024.

\bibitem[Fleiss(1971)]{fleiss1971}
Fleiss, J.~L.
\newblock Measuring nominal scale agreement among many raters.
\newblock \emph{Psychological Bulletin}, 76(5):378--382, 1971.

\bibitem[Chen et al.(2023)]{chen2023usc}
Chen, X., Aksitov, R., Alon, U., Ren, J., Xiao, K., Yin, P., Prakash, S., Sutton, C., Wang, X., and Zhou, D.
\newblock Universal self-consistency for large language model generation.
\newblock \emph{arXiv:2311.17311}, 2023.

\bibitem[Li et al.(2023)]{li2023diverse}
Li, Y., Lin, Z., Zhang, S., Fu, Q., Chen, B., Lou, J.-G., and Chen, W.
\newblock Making language models better reasoners with step-aware verifier.
\newblock In \emph{ACL}, 2023.

\bibitem[Taubenfeld et al.(2025)]{taubenfeld2025weighted}
Taubenfeld, A., et al.
\newblock Weighted self-consistency for language model reasoning.
\newblock 2025.

\bibitem[Loo et al.(2025)]{loo2025sc}
Loo, N., et al.
\newblock On the scaling limits of self-consistency.
\newblock 2025.

\bibitem[{DyLoS}(2026)]{dylos2026}
{DyLoS Authors}.
\newblock {DyLoS}: Dynamic logic solving with language models.
\newblock In \emph{EACL}, 2026.

\bibitem[{Aristotle}(2025)]{aristotle2025}
{Aristotle Authors}.
\newblock Aristotle: Structured reasoning with external verification.
\newblock In \emph{ACL}, 2025.

\bibitem[{FoVer}(2025)]{fover2025}
{FoVer Authors}.
\newblock {FoVer}: Formal verification of large language model outputs.
\newblock \emph{TACL}, 2025.

\bibitem[Snell et al.(2025)]{snell2025scaling}
Snell, C., et al.
\newblock Scaling {LLM} test-time compute.
\newblock 2025.

\bibitem[{DeepSeek-AI}(2025)]{deepseekr1}
{DeepSeek-AI}.
\newblock {DeepSeek-R1}: Incentivizing reasoning capability in {LLMs} via reinforcement learning.
\newblock \emph{arXiv:2501.12948}, 2025.

\bibitem[Cobbe et al.(2021)]{cobbe2021orm}
Cobbe, K., Kosaraju, V., Bavarian, M., Chen, M., Jun, H., Kaiser, L., Plappert, M., Tworek, J., Hilton, J., Nakano, R., Hesse, C., and Schulman, J.
\newblock Training verifiers to solve math word problems.
\newblock \emph{arXiv:2110.14168}, 2021.

\bibitem[Cobbe et al.(2021b)]{cobbe2021gsm8k}
Cobbe, K., Kosaraju, V., Bavarian, M., Chen, M., Jun, H., Kaiser, L., Plappert, M., Tworek, J., Hilton, J., Nakano, R., Hesse, C., and Schulman, J.
\newblock Training verifiers to solve math word problems.
\newblock \emph{arXiv:2110.14168}, 2021.

\bibitem[Lightman et al.(2023)]{lightman2023prm}
Lightman, H., Kosaraju, V., Burda, Y., Edwards, H., Baker, B., Lee, T., Leike, J., Schulman, J., Sutskever, I., and Alignment, K.
\newblock Let's verify step by step.
\newblock In \emph{ICLR}, 2024.

\bibitem[Liang et al.(2023)]{liang2023debate}
Liang, T., He, Z., Jiao, W., Wang, X., Wang, Y., Wang, R., Yang, Y., Tu, Z., and Shi, S.
\newblock Encouraging divergent thinking in large language models through multi-agent debate.
\newblock \emph{arXiv:2305.19118}, 2023.

\bibitem[Du et al.(2023)]{du2023debate}
Du, Y., Li, S., Torralba, A., Tenenbaum, J.~B., and Mordatch, I.
\newblock Improving factuality and reasoning in language models through multiagent debate.
\newblock \emph{arXiv:2305.14325}, 2023.

\bibitem[Wei et al.(2022)]{wei2022cot}
Wei, J., Wang, X., Schuurmans, D., Bosma, M., Ichter, B., Xia, F., Chi, E., Le, Q., and Zhou, D.
\newblock Chain-of-thought prompting elicits reasoning in large language models.
\newblock In \emph{NeurIPS}, 2022.

\bibitem[Kojima et al.(2022)]{kojima2022zerocot}
Kojima, T., Gu, S.~S., Reid, M., Matsuo, Y., and Iwasawa, Y.
\newblock Large language models are zero-shot reasoners.
\newblock In \emph{NeurIPS}, 2022.

\bibitem[Han et al.(2022)]{han2022folio}
Han, S., Schoelkopf, H., Zhao, Y., Qi, Z., Riddell, M., Benson, L., Sun, L., Zubova, E., Qiao, Y., Burtell, M., et al.
\newblock {FOLIO}: Natural language reasoning with first-order logic.
\newblock \emph{arXiv:2209.00840}, 2022.

\bibitem[Tafjord et al.(2021)]{tafjord2021proofwriter}
Tafjord, O., Dalvi, B., and Clark, P.
\newblock {ProofWriter}: Generating implications, proofs, and abductive statements over natural language.
\newblock In \emph{Findings of ACL}, 2021.

\bibitem[Clark et al.(2018)]{clark2018arc}
Clark, P., Cowhey, I., Etzioni, O., Khot, T., Sabharwal, A., Schoenick, C., and Tafjord, O.
\newblock Think you have solved question answering? {T}ry {ARC}, the {AI2} reasoning challenge.
\newblock \emph{arXiv:1803.05457}, 2018.

\bibitem[Geva et al.(2021)]{geva2021strategyqa}
Geva, M., Khashabi, D., Segal, E., Khot, T., Roth, D., and Berant, J.
\newblock Did {A}ristotle use a laptop? {A} question answering benchmark with implicit reasoning strategies.
\newblock \emph{TACL}, 2021.

\bibitem[Liu et al.(2023)]{liu2023logiqa}
Liu, J., et al.
\newblock {LogiQA} 2.0: An improved dataset for logical reasoning.
\newblock \emph{IEEE/ACM Transactions on Audio, Speech, and Language Processing}, 2023.

\bibitem[Jiang et al.(2023)]{jiang2023mistral}
Jiang, A.~Q., Sablayrolles, A., Mensch, A., Bamford, C., Chaplot, D.~S., de~las Casas, D., Bressand, F., Lengyel, G., Lample, G., et al.
\newblock Mistral 7{B}.
\newblock \emph{arXiv:2310.06825}, 2023.

\bibitem[{Qwen Team}(2024)]{qwen2024}
{Qwen Team}.
\newblock Qwen2.5 technical report.
\newblock 2024.

\bibitem[Touvron et al.(2023)]{touvron2023llama}
Touvron, H., Martin, L., Stone, K., Albert, P., Almahairi, A., et al.
\newblock {LLaMA} 2: Open foundation and fine-tuned chat models.
\newblock \emph{arXiv:2307.09288}, 2023.

\bibitem[{SATBench}(2025)]{satbench2025}
{SATBench Authors}.
\newblock {SATBench}: Benchmarking constraint satisfaction with language models.
\newblock In \emph{EMNLP}, 2025.

\bibitem[{LLM-CP}(2025)]{llmcp2025}
{LLM-CP Authors}.
\newblock {LLM} meets constraint propagation.
\newblock In \emph{IJCAI}, 2025.

\end{thebibliography}

\end{document}
